\section{Related Work}

Agent decision-making and task planning lie at the core of both single-agent and multi-agent success. Foundational cognitive architectures such as CoALA~\cite{sumers2024cognitivearchitectures} and ReAct~\cite{yao2023react} underpin most modern agent implementations but assumed relatively small external tool sets. As agents face real-world tasks, this assumption breaks down: long contexts cause information loss~\cite{liu2023lostmiddlellms} and task distraction~\cite{shi2023llmseasily}, incorrect tool Retrieval accounts for over half of task failures~\cite{mo2026livemcpbench}, and even state-of-the-art models suffer significant performance limitations without dedicated retrieval optimizations~\cite{luo2025mcpuniverse}. This has driven a growing body of work on effective tool selection from large and evolving resource sets~\cite{patil2023gorilla, qin2023toolllm, gaurav2025dynamicreact, wang2025toolgen}.

Much of this literature targets MCP, which established itself early as a widely adopted tool-integration standard~\cite{ehtesham2025surveyagentinteroperabilityprotocols}. Approaches to improving MCP tool discovery can thereby be divided into LLM-internal and external retrieval solutions. Internal approaches such as ToolGen~\cite{wang2025toolgen} represent each tool as a token and fine-tune the model to retrieve them. While effective even at scale, this requires fixed token-level representations and explicit fine-tuning constraints ill-suited to rapidly evolving tool ecosystems.

External approaches decouple retrieval from the model and can be categorized by strategy. Proxy-based solutions such as MCP Bridge or ToolRegistry~\cite{ahmadi2026mcpbridge, ding2025toolregistry} implement a lightweight gateway mediating communication on behalf of the agent. Semantic retrieval approaches~\cite{gaurav2025dynamicreact, gan2025ragmcp, fei2025mcpzero, lumer2025tooltoagent, lumer2025scalemcp, gaurav2025codemem} introduce vector-based selection of relevant tools at query time and differ primarily in their embedding strategies, retrieval metrics, or specialized description and execution formats. Lastly, graph-based approaches~\cite{nizar2025agentasagraph, nie2025llmtoolretrieval} encode agent and tool relationships in knowledge graphs to discover previously unseen connections, but require explicit graph construction and assume agents can freely share their tools, which rarely holds in industry settings with independent peers.

Collectively, this literature reflects a shift toward agents that dynamically fetch resources based on semantic relevance rather than passively receiving static capability dumps. This resonates with our core argument that agents should govern their own capability acquisition, progressively gathering only the information relevant to their current task instead of being overloaded upfront. 

Yet most work remains focused on tool calling alone. Real-world settings, however, also require inter-agent communication. The proliferation of agent-communication protocols~\cite{ehtesham2025surveyagentinteroperabilityprotocols}, most notably A2A, compounds the existing challenges: Context bloat, protocol confusion, and retrieval errors all worsen when multiple protocols must coexist. 
Our work addresses this gap by proposing a gateway that abstracts both tools and peer agents into a unified capability model, hiding protocol-specific details behind a common interface. This design not only scales more gracefully across protocols than direct exposure but also integrates semantic retrieval at the gateway level, combining insights from proxy-based mediation and tool-retrieval research into a single architecture.